% Section 7: Spectral Analysis of Linear Transformations  (original handwritten pages 52--58)
\section{Spectral Analysis of Linear Transformations}\label{sec:spectral}
\syllabusweek{Weeks 8--9 --- eigenvalues and eigenvectors, invariance under operators, block diagonal forms}
\subsection{Invariant Subspaces}\label{ss:invariant}
For a matrix $A$, the set of its eigenvalues $\{\lambda_1,\lambda_2,\dots,\lambda_n\}$ is called the \term{spectrum} of the matrix.
\begin{definition}
Let $\mathcal{A}:\V\to\V$ be a linear transformation defined over a vector space $\V$. A subspace $M$ of $\V$ is said to be \term{invariant under $\mathcal{A}$} iff
\[
\mathcal{A}x\in M\quad\forall x\in M
\]
\end{definition}
\begin{center}
\begin{tikzpicture}
\draw (0,0) ellipse (0.8 and 0.9);
\draw (0.05,0.05) circle (0.42);
\node[font=\scriptsize] at (-0.45,0.62) {$M$};
\node[font=\scriptsize] at (0,-1.15) {$\V$};
\draw (3.2,0.1) ellipse (0.7 and 0.9);
\draw (3.2,0.15) circle (0.42);
\node[font=\scriptsize] at (3.2,0.78) {$M$};
\node[font=\scriptsize] at (3.85,-0.85) {$\V$};
\fill (-0.1,-0.05) circle (1pt);
\fill (3.1,0.2) circle (1pt);
\draw[-{Stealth[length=1.5mm]}] (-0.1,-0.05) -- node[above]{$\mathcal{A}$} (3.05,0.2);
\end{tikzpicture}
\end{center}

\begin{example}
$\V=\{f\mid f\text{ functions defined over }\R\text{ which are infinitely differentiable}\}$
\[
\mathcal{A}=\frac{d}{dt}:\ \V\to\V, \qquad M_1\triangleq\Span\{\sin t\}\subset\V
\]
\[
\mathcal{A}(\alpha\sin t)=\alpha\mathcal{A}\sin t=\alpha\cos t\notin M_1
\]
$\therefore$ \key{$M_1$ is not invariant under $\mathcal{A}$.}
\begin{align*}
M_2&\triangleq\Span\{\cos t,\sin t\}\\
x&=\alpha_1\cos t+\alpha_2\sin t\\
\mathcal{A}x&=-\alpha_1\sin t+\alpha_2\cos t\in\Span\{\cos t,\sin t\}= M_2
\end{align*}
$\therefore$ \key{$M_2$ is invariant under $\mathcal{A}$.}
\begin{align*}
M_3&=\{0\}, \qquad \mathcal{A}0=0\in M_3\\
M_4&=\Span\{e^{2t}\},\qquad x=\alpha e^{2t},\qquad \mathcal{A}x=2\alpha e^{2t}\in M_4
\end{align*}
$\therefore$ $M_3$ and $M_4$ are invariant under $\mathcal{A}$ ($e^{2t}$ is mapped to a multiple of itself: it is an ``eigenfunction'' of $\frac{d}{dt}$ with eigenvalue $2$).
\end{example}
% corrected: the notes write "not in M" for M_1; made explicit as M_1.
\begin{insight}
The trivial subspaces $\{0\}$ and $\V$ are always invariant. \key{A one-dimensional subspace $\Span\{e\}$ is invariant iff $\mathcal{A}e=\lambda e$} for some scalar $\lambda$, i.e.\ iff $e$ is an eigenvector --- this is the link between invariant subspaces and the eigenvalues defined below.
\end{insight}

% ===== Page 52 =====
\origpage{52}
\needspace{8\baselineskip}
\begin{example}
$\Null(\mathcal{A})$ \& $\Range(\mathcal{A})$ are invariant under $\mathcal{A}$.

\begin{minipage}[t]{0.48\textwidth}
\paragraph{$\Null(\mathcal{A})$.} $x\in\Null(\mathcal{A})$
\[
\mathcal{A}x=0_{\V}\in\Null(\mathcal{A})
\]
$\Null(\mathcal{A})$ is invariant under $\mathcal{A}$.
\end{minipage}\hfill
\begin{minipage}[t]{0.48\textwidth}
\paragraph{$\Range(\mathcal{A})$.} $x\in\Range(\mathcal{A})$
\[
\mathcal{A}x\in\Range(\mathcal{A})
\]
$\Range(\mathcal{A})$ is invariant under $\mathcal{A}$.
\end{minipage}
\end{example}

\subsection{Powers and Polynomials of a Linear Transformation}\label{ss:powers}
\begin{definition}
\term{Powers} of a linear transformation $\mathcal{A}:\V\to\V$ are defined as
\[
\mathcal{A}^k x=\underbrace{\mathcal{A}(\mathcal{A}\cdots\mathcal{A}(\mathcal{A}x)\dots)}_{\mathcal{A}\text{ is applied }k\text{ times}}
\]
\end{definition}
When the domain and the codomain are the same, we can define the powers of a linear transformation.
\begin{itemize}[nosep,leftmargin=1.5em]
\item $\R^n\to\R^m$, $m\neq n$ $\Rightarrow$ not defined.
\item $\mathcal{A}:\V\to\W$: \ $\mathcal{A}(\underbrace{\mathcal{A}x}_{\in\W})$ \ $\times$ not defined, since $\mathcal{A}$ acts on elements $\in\V$.
\end{itemize}

\sep
\point \term{Polynomials of $\mathcal{A}$} can be constructed using linear combinations of powers of $\mathcal{A}$.
\begin{center}
\begin{minipage}{0.66\textwidth}
\centering$\displaystyle\begin{aligned}
p(s)&=\alpha_0s^n+\alpha_1s^{n-1}+\dots+\alpha_{n-1}s+\alpha_n\,(s^0)\\
p(\mathcal{A})&=\alpha_0\mathcal{A}^n+\alpha_1\mathcal{A}^{n-1}+\dots+\alpha_{n-1}\mathcal{A}^1+\alpha_n\mathcal{A}^0\\
p(\mathcal{A})&=\alpha_0\mathcal{A}^n+\dots+\alpha_{n-1}\mathcal{A}+\alpha_nI
\end{aligned}$
\end{minipage}\hfill
\begin{minipage}{0.28\textwidth}
\begin{keybox}
\centering
$\mathcal{A}^0x=x=Ix$\\
$\mathcal{A}^0=I$
\end{keybox}
\end{minipage}
\end{center}
For example, with $p(\mathcal{A})=3\mathcal{A}^3+2\mathcal{A}^2+\mathcal{A}+I$:
\begin{align*}
p(\mathcal{A})\mathcal{A}x&=(3\mathcal{A}^3+2\mathcal{A}^2+\mathcal{A}+I)\mathcal{A}x=3\mathcal{A}^4x+2\mathcal{A}^3x+\mathcal{A}^2x+\mathcal{A}x\\
&=\mathcal{A}(3\mathcal{A}^3x+2\mathcal{A}^2x+\mathcal{A}x+x)\\
&=\mathcal{A}\underbrace{(3\mathcal{A}^3+2\mathcal{A}^2+\mathcal{A}+I)}_{p(\mathcal{A})}x
\end{align*}
\[
\boxed{p(\mathcal{A})\mathcal{A}x=\mathcal{A}p(\mathcal{A})x\quad\forall x\in\V}
\qquad\Rightarrow\qquad
p(\mathcal{A})\mathcal{A}=\mathcal{A}p(\mathcal{A})
\]
\explain{\key{$\mathcal{A}$ commutes with $p(\mathcal{A})$.}}

\begin{example}
$\Null(p(\mathcal{A}))$ \& $\Range(p(\mathcal{A}))$ are invariant under $\mathcal{A}$.

\paragraph{$\Null(p(\mathcal{A}))$.} $x\in\Null(p(\mathcal{A}))$. \ Is $\mathcal{A}x\in\Null(p(\mathcal{A}))$?
\begin{align*}
p(\mathcal{A})\mathcal{A}x&=\mathcal{A}\underbrace{p(\mathcal{A})x}_{0_{\V}}=\mathcal{A}0_{\V}=0_{\V}\\
&\Rightarrow\mathcal{A}x\in\Null(p(\mathcal{A}))
\end{align*}

% ===== Page 53 =====
\needspace{6\baselineskip}
\origpage{53}%
\paragraph{$\Range(p(\mathcal{A}))$.} $x\in\Range(p(\mathcal{A}))$, \ $x=p(\mathcal{A})y$.
\begin{align*}
\mathcal{A}x&=\mathcal{A}p(\mathcal{A})y=p(\mathcal{A})\underbrace{\mathcal{A}y}_{y'}=p(\mathcal{A})y'\\
&\Rightarrow\mathcal{A}x\in\Range(p(\mathcal{A}))
\end{align*}
\end{example}

\subsection{Eigenvalues and Eigenvectors}\label{ss:eigen}
\point Matrix representation:
\[
\mathcal{A}:\underset{\substack{\downarrow\\ B\\ \text{basis set}\\ \text{for }\V}}{\overset{n}{\V}}\to\underset{\substack{\downarrow\\ B\\ \text{basis set}\\ \text{for }\V}}{\overset{n}{\V}}
\qquad A:\C^n\to\C^n
\qquad
\begin{aligned}
y&=\mathcal{A}x &\quad \lambda e&=\mathcal{A}e\\
&\Downarrow & &\Downarrow\\
[y]_B&=A[x]_B & \lambda[e]_B&=A[e]_B
\end{aligned}
\]
\[
\det\mathcal{A}\triangleq\det A,\qquad \tr\mathcal{A}\triangleq\tr A
\]
% corrected: notes had "tr A := det A" (slip; the trace of the operator is the trace of its matrix)
\explain{Well defined: another basis gives $\bar A=P^{-1}AP$, and $\det(P^{-1}AP)=\det A$, $\tr(P^{-1}AP)=\tr A$.}

\begin{definition}
Let $A$ denote the matrix representation of a linear transformation $\mathcal{A}$ from $\V$ into $\V$ ($A$ is a square matrix). The \term{eigenvalues} of $\mathcal{A}$ denoted by $\lambda_i$ are defined as the $n$ roots of the equation
\[
\det(sI-A)=0,
\]
where $d(s)\triangleq\det(sI-A)$ is called the \term{characteristic polynomial} of $A$.
The vectors $e_i\neq0$ satisfying $Ae_i=\lambda_i e_i$ are called \term{eigenvectors} of $A$.
\end{definition}

\begin{center}
\begin{minipage}{0.48\textwidth}
\centering
\begin{tikzpicture}[scale=0.75,>={Stealth[length=1.8mm]}]
\draw (-3,0)--(3,0);
\draw[->] (0,-1.5)--(0,2.8);
\draw[->] (0,0)--(0.25,2.6) node[above right]{$\mathcal{A}x$};
\draw[->] (0,0)--(0.8,1.9) node[above]{$x$};
\draw[->] (0,0)--(2.6,1.3) node[above]{$2e_i$};
\draw[->,very thick] (0,0)--(1.3,0.65) node[below right]{$e_i$};
\draw[->] (0,0)--(-2.6,-1.3) node[below]{$-2e_i$};
\end{tikzpicture}
\end{minipage}\hfill
\begin{minipage}{0.45\textwidth}
\begin{align*}
&\mathcal{A}e_i=\lambda_ie_i\\
&\lambda_i=2\\
&\lambda_i=-2
\end{align*}
\explain{A general $x$ is turned into a different direction ($\mathcal{A}x\not\parallel x$); an eigenvector $e_i$ is only scaled, e.g.\ to $2e_i$ if $\lambda_i=2$ or to $-2e_i$ if $\lambda_i=-2$, so the line $\Span\{e_i\}$ is invariant.}
\end{minipage}
\end{center}

\sep
\[
A=\underset{\text{rotation matrix}}{\bmat{0&-1\\1&0}}\qquad
x_1=\bmat{1\\0}\qquad x_2=\bmat{0\\1}\qquad x_3=\bmat{1\\1}
\]

\begin{minipage}{0.42\textwidth}
\centering
\begin{tikzpicture}[scale=1.1,>={Stealth[length=1.8mm]}]
\draw[->] (-1.7,0)--(1.8,0);
\draw[->] (0,-1.0)--(0,1.6);
\draw[->,thick] (0,0)--(1,0) node[below]{$x_1$};
\draw[->] (-0.04,0)--(-0.04,1) node[above left,inner sep=1pt]{$x_2$};
\draw[->,very thick] (0.03,0)--(0.03,1) node[above right,inner sep=1pt]{$\mathcal{A}x_1$};
\draw[->] (0,0.03)--(-1,0.03) node[below]{$\mathcal{A}x_2$};
\draw[->] (0,0)--(1,1) node[right]{$x_3$};
\draw[->] (0,0)--(-1,1) node[left]{$\mathcal{A}x_3$};
\end{tikzpicture}
\end{minipage}\hfill
\begin{minipage}{0.55\textwidth}
\begin{align*}
&s^2+1=0\qquad\lambda_{1,2}=\pm j\\
&\underset{\text{real}}{A}e_i=\underset{\text{complex}}{\lambda_i}e_i\qquad R:\text{ rotation matrix}\\
&Re_i=\lambda_ie_i\\
&\norm{Re_i}=\norm{\lambda_ie_i}\\
&\cancel{\norm{e_i}}=\abs{\lambda_i}\cancel{\norm{e_i}}\qquad (e_i:\text{ nonzero vector})\\
&\abs{\lambda_i}=1\\
&\therefore\ \boxed{\lambda_i=e^{j\omega}}
\end{align*}
\end{minipage}
% corrected: notes had "R: rotation vectors"; R denotes the rotation matrix
\needspace{7\baselineskip}
\begin{insight}
No real vector keeps its direction under a $90^\circ$ rotation, so there are no real eigenvectors; the eigenvectors are complex: $A[1\ {-j}]\T=[j\ 1]\T=j\,[1\ {-j}]\T$ and $A[1\ j]\T=-j\,[1\ j]\T$. A rotation preserves length, $\norm{Re}=\norm{e}$ (also for complex $e$, since $R\Hm R=I$), which is what forces $\abs{\lambda_i}=1$; here $\lambda=\pm j=e^{\pm j\pi/2}$, i.e.\ $\omega=\pm\pi/2$ is the rotation angle.
\end{insight}

% ===== Page 54 =====
\origpage{54}
\sep
\point Existence of eigenvectors:
\[
A\overset{\neq0}{e_i}=\lambda_ie_i\qquad d(\lambda_i)=\det(\lambda_iI-A)=0
\]
\begin{align*}
&A=[c_1\ c_2\ \dots\ c_n]\\
&\exists\ \alpha_1,\alpha_2,\dots,\alpha_n\text{ not all equal to zero satisfying}\\
&\qquad\alpha_1c_1+\alpha_2c_2+\dots+\alpha_nc_n=0\\
&\qquad\underset{A}{[c_1\ c_2\ \dots\ c_n]}\bmat{\alpha_1\\\alpha_2\\\vdots\\\alpha_n}=0
\end{align*}
\explain{$A$ is not full column rank $\Rightarrow$ $\exists\alpha\neq0$ such that $A\alpha=0$. Apply this to the singular matrix $\lambda_iI-A$ ($\det(\lambda_iI-A)=0$ means its columns are dependent):}
\begin{align*}
&\exists\alpha\ \text{s.t.}\ (\lambda_iI-A)\alpha=0\\
&\qquad A\alpha=\lambda_i\alpha\ \Rightarrow\ \text{Eğer bir eigenvalue varsa en az 1 eigenvector vardır}
\end{align*}
\explain{\key{If there is an eigenvalue, then there is at least 1 eigenvector}: $\alpha\neq0$ is an eigenvector for $\lambda_i$.}

\begin{example}
$A=\bmat{2&1\\1&2}$:
\begin{align*}
d(s)&=\det(sI-A)=\det\bmat{s-2&-1\\-1&s-2}\\
&=(s-2)^2-1=s^2-4s+3=(s-3)(s-1)=0\\
&\Rightarrow\ \lambda_1=3,\quad\lambda_2=1
\end{align*}
\begin{align*}
&Ae_i=\lambda_ie_i\\
&(A-\lambda_iI)e_i=0\\
&e_i\in\underbrace{\Null(A-\lambda_iI)}_{\text{\term{eigenspace}}},\quad\dim\geq1
\qquad\qquad
\begin{aligned}
p(s)&=s-\lambda_i\\
p(A)&=A-\lambda_iI
\end{aligned}
\end{align*}
\begin{alignat*}{2}
A-\lambda_1I&=\bmat{-1&1\\1&-1},&\qquad e_1&=\bmat{1\\1}\\
A-\lambda_2I&=\bmat{1&1\\1&1},&\qquad e_2&=\bmat{1\\-1}
\end{alignat*}
\begin{align*}
\Null(A-\lambda_1I)&=\Span\left\{\bmat{1\\1}\right\}\\
\Null(A-\lambda_2I)&=\Span\left\{\bmat{1\\-1}\right\}
\end{align*}
\[
\Null(A-\lambda_1I)\overset{\perp}{\oplus}\Null(A-\lambda_2I)=\C^2
\]
% corrected/simplified: notes put perp marks over each N(.) and "dim = 2" marks over C^2
\explain{$\overset{\perp}{\oplus}$: basis vectors are linearly indep.\ and orthogonal ($[1\ 1]\T\perp[1\ {-1}]\T$); dimensions $1+1=2=\dim\C^2$.}
\[
\alpha\bmat{1\\1}+\beta\bmat{1\\-1}=\begin{matrix}\R^2\\\C^2\end{matrix}
\]
\[
\alpha\bmat{1\\1}+\beta\bmat{1\\-1}=\bmat{0\\0}\quad\text{iff }\alpha=\beta=0
\ \Rightarrow\ \text{subspaces are linearly indep.}
\]
\explain{Orthogonality is special here because $A=A\Hm$ is Hermitian; in general eigenspaces are only linearly independent (see p.~57).}
\end{example}

% ===== Page 55 =====
\origpage{55}
\sep
\point Change of basis:
\[
y=Ax,\qquad B=\left\{\bmat{1\\0},\bmat{0\\1}\right\}\ \to\ \text{canonical basis set}
\]
\[
x=\bmat{x_1\\x_2},\quad [x]_B,\qquad [y]_B=A[x]_B\ \sim\ [y]_{\bar{B}}=\bar{A}[x]_{\bar{B}},
\qquad
\bmat{x_1\\x_2}=x_1\bmat{1\\0}+x_2\bmat{0\\1}
\]
\[
\bar{B}=\left\{\bmat{1\\1},\bmat{1\\-1}\right\},\qquad A=\bmat{2&1\\1&2}
\]
\begin{align*}
x&=\bmat{1\\1}, & [x]_{\bar{B}}&=\bmat{1\\0}, & y=Ax&=\bmat{3\\3}=3\bmat{1\\1}, & [y]_{\bar{B}}&=\bmat{3\\0}\\
&& \bmat{3\\0}&=\bar{A}\bmat{1\\0} & \Rightarrow\ \bar{A}&=\bmat{3&\\0&}\\[4pt]
x&=\bmat{1\\-1}, & [x]_{\bar{B}}&=\bmat{0\\1}, & y=Ax&=\bmat{1\\-1}=1\bmat{1\\-1}, & [y]_{\bar{B}}&=\bmat{0\\1}\\
&& \bmat{0\\1}&=\bar{A}\bmat{0\\1} & \Rightarrow\ \bar{A}&=\bmat{3&0\\0&1}
\end{align*}

\sep
\[
P=[e_1\mid e_2],\qquad \bar{A}=f(A)=\ ?
\]
\[
\begin{aligned}Ae_1&=\lambda_1e_1\\Ae_2&=\lambda_2e_2\end{aligned}
\]
\needspace{9\baselineskip}% keep the derivation together with its explanation below
\begin{align*}
[\lambda_1e_1\mid\lambda_2e_2]&=[Ae_1\mid Ae_2]=A[e_1\mid e_2]\\
\underbrace{[e_1\mid e_2]}_{P}\underbrace{\bmat{\lambda_1&0\\0&\lambda_2}}_{\bar{A}}&=A\cdot P
\ \Rightarrow\ P\bar{A}=AP\\
&\Rightarrow\ \boxed{\bar{A}=P^{-1}AP}
\end{align*}
\explain{$P$ is invertible (full rank: its columns $e_1,e_2$ are linearly indep.).}
\begin{insight}
The columns of $P$ are the new basis vectors written in the old basis, so $[x]_B=P[x]_{\bar B}$. Then $[y]_B=A[x]_B$ becomes $P[y]_{\bar B}=AP[x]_{\bar B}$, i.e.\ $[y]_{\bar B}=P^{-1}AP\,[x]_{\bar B}$. Check:
\[
P=\bmat{1&1\\1&-1},\qquad P^{-1}=\tfrac12\bmat{1&1\\1&-1},\qquad
P^{-1}AP=\tfrac12\bmat{1&1\\1&-1}\bmat{3&1\\3&-1}=\bmat{3&0\\0&1},
\]
as found above.
\end{insight}
\[
\C^2=\underset{1}{\Null(A-\lambda_1I)}\overset{\perp}{\oplus}\underset{1}{\Null(A-\lambda_2I)},
\qquad
\left[\begin{array}{c|c}\overset{1\times1}{\lambda_1}&0\\\hline 0&\lambda_2\end{array}\right]
\]
\explain{Each block is $1\times1$ as $\dim\Null(A-\lambda_1I)=\dim\Null(A-\lambda_2I)=1$.}

\subsection{Block Diagonal Representations}\label{ss:block-diag}
\begin{theorem}[Block diagonal representation]
Consider the linear transformation
\[
y=Ax\quad\text{where }A\in\C^{n\times n},\ y,x\in\C^n.
\]
Suppose that
\begin{enumerate}[label=\roman*)]
\item $\C^n=M_1\oplus M_2\oplus\dots\oplus M_k$.
\item $M_i$ is invariant under $A$ $\forall i$.
\end{enumerate}
Then, the transformation $A$ has a \key{block diagonal representation}
\[
\bar{A}=\bmat{\bar{A}_1&&&0\\&\bar{A}_2&&\\&&\ddots&\\0&&&\bar{A}_k}
\]
% ===== Page 56 =====
\origpage{56}%
where $\bar{A}=P^{-1}AP$,
\begin{align*}
P&=[P_1\mid P_2\mid\dots\mid P_k]\\
P_i&=[b_i^1\ b_i^2\ b_i^3\ \dots\ b_i^{n_i}],\qquad n_i=\dim M_i
\end{align*}
$b_i^1,\dots,b_i^{n_i}\in\C^n$ are basis vectors of $M_i$. $\bar{A}_i$ is of size $n_i\times n_i$.
\end{theorem}

\begin{proof}[Proof sketch]
\begin{align*}
&k=2,\qquad \C^n=M_1\oplus M_2,\qquad n=5\\
&M_1=\Span\{b_1^1,b_1^2\},\qquad\dim M_1=2\\
&M_2=\Span\{b_2^1,b_2^2,b_2^3\},\qquad\dim M_2=3
\end{align*}
% corrected: notes had M_1 = sp{b_1^1, b_2^2}, M_2 = sp{b_1^1, b_2^2, b_3^3}; indices fixed to match P_1, P_2 below (b_i^j: j-th basis vector of M_i)
\[
P=[P_1\mid P_2],\qquad
\begin{aligned}
P_1&=[b_1^1\ b_1^2]_{5\times2}\\
P_2&=[b_2^1\ b_2^2\ b_2^3]_{5\times3}
\end{aligned}
\qquad
\bar{A}=P^{-1}AP\ \Rightarrow\ P\bar{A}=AP
\]
\begin{align*}
\underbrace{[b_1^1\,b_1^2\mid b_2^1\,b_2^2\,b_2^3]}_{P}
\underbrace{\overset{2\times2\ \to\ \dim M_1\times\dim M_1}{\left[\begin{array}{cc|ccc}
\alpha_1^1&\beta_1^1&0&0&0\\
\alpha_1^2&\beta_1^2&0&0&0\\\hline
0&0&\alpha_2^1&\beta_2^1&\gamma_2^1\\
0&0&\alpha_2^2&\beta_2^2&\gamma_2^2\\
0&0&\alpha_2^3&\beta_2^3&\gamma_2^3
\end{array}\right]}}_{\bar{A}\quad\left(\substack{3\times3\\ \dim M_2\times\dim M_2}\right)}
&=A\overbrace{[b_1^1\ b_1^2\ b_2^1\ b_2^2\ b_2^3]}^{P}\\
&=[\underbrace{Ab_1^1}_{\in M_1}\mid\underbrace{Ab_1^2}_{\in M_1}\mid\underbrace{Ab_2^1}_{\in M_2}\mid\underbrace{Ab_2^2}_{\in M_2}\mid\underbrace{Ab_2^3}_{\in M_2}]\qedhere
\end{align*}
\end{proof}
\explain{$\in M_1$: lin.\ comb.\ of $M_1$ vectors; \ $\in M_2$: lin.\ comb.\ of $M_2$ vectors.}
\begin{insight}[Why the off-diagonal blocks are zero]
Column $j$ of $P\bar A$ is $P$ times column $j$ of $\bar A$, i.e.\ the combination of the basis vectors with the coefficients in that column. Since $M_1$ is invariant, $Ab_1^1=\alpha_1^1b_1^1+\alpha_1^2b_1^2$ needs \emph{no} $M_2$ vectors, so the lower entries of columns 1--2 are zero; likewise $Ab_2^j\in M_2$ gives zeros in the upper part of columns 3--5. The coefficients are unique because the five vectors form a basis of $\C^5$ ($\C^n=M_1\oplus M_2$), which also makes $P$ invertible. Each block $\bar A_i$ is the matrix of $A$ restricted to $M_i$.
\end{insight}

\begin{example}
\[
A=\bmat{1&1&-1\\-1&3&-2\\0&0&1}
\]
\[
M_1\triangleq\Span\left\{\overset{b_1^1}{\bmat{1\\1\\0}},\overset{b_1^2}{\bmat{0\\1\\0}}\right\},\ \dim M_1=2,
\qquad
M_2\triangleq\Span\left\{\overset{b_2^1}{\bmat{0\\1\\1}}\right\},\ \dim M_2=1
\]
% corrected: notes had M_2 = { b_2^1 } (set instead of span)
\begin{align*}
Ab_1^1&=\bmat{2\\2\\0}=2b_1^1+0b_1^2\in M_1\\
Ab_1^2&=\bmat{1\\3\\0}=1b_1^1+2b_1^2\in M_1\\
A(\alpha_1b_1^1+\alpha_2b_1^2)&=\alpha_1Ab_1^1+\alpha_2Ab_1^2\\
&=\alpha_1(2b_1^1)+\alpha_2(1b_1^1+2b_1^2)\\
&=(2\alpha_1+\alpha_2)b_1^1+2\alpha_2b_1^2\in M_1
\end{align*}
$\therefore$ \key{$M_1$ is invariant under $A$.}
\[
\left.\begin{aligned}
Ab_1^1&=2b_1^1 &&\Rightarrow\ b_1^1\text{ is an eigenvector}\\
Ab_1^2&=1b_1^1+2b_1^2 &&\Rightarrow\ b_1^2\text{ is not an eigenvector}
\end{aligned}\right\}\ \Rightarrow\ M_1\text{ is not an eigenspace}
\]
\[
Ab_2^1=\bmat{0\\1\\1}=1b_2^1\in M_2\ \Rightarrow\
\begin{aligned}
&\therefore\ M_2\text{ is invariant under }A\\
&M_2\text{ is an eigenspace}
\end{aligned}
\]

% ===== Page 57 =====
\origpage{57}
\[
\C^3=\underset{\checkmark}{M_1}\oplus M_2,\qquad
\alpha_1b_1^1+\alpha_2b_1^2+\beta_1b_2^1=0
\]
\explain{Linearly independent (only solution $\alpha_1=\alpha_2=\beta_1=0$: the 3rd entry gives $\beta_1=0$, the 1st $\alpha_1=0$, then $\alpha_2=0$). No orthogonality (no need for it actually).}
\[
\bar{A}=\left[\begin{array}{cc|c}\ \ &\ \ &\ \ \\ \ &\ &\ \\\hline\ &\ &\ \end{array}\right],
\qquad
\bar{A}=P^{-1}AP\Rightarrow P\bar{A}=AP,
\qquad
P=[b_1^1\ b_1^2\mid b_2^1]
\]
% corrected: notes had P = [b_1^1 b_2^1 | b_2^1]; the basis of M_1 followed by that of M_2 is [b_1^1 b_1^2 | b_2^1], as in the next line.
\begin{align*}
\underbrace{[b_1^1\ b_1^2\ b_2^1]}_{P}
\underbrace{\left[\begin{array}{cc|c}2&1&0\\0&2&0\\\hline0&0&1\end{array}\right]}_{\bar{A}}
&=A[b_1^1\ b_1^2\ b_2^1]\\
&=[Ab_1^1\mid Ab_1^2\mid Ab_2^1]
\end{align*}
\explain{The columns of $\bar A$ are the coefficients found above: $Ab_1^1=2b_1^1$, $Ab_1^2=1b_1^1+2b_1^2$, $Ab_2^1=1b_2^1$.}
\[
d(s)=\det(sI-A)=(s-2)^{\text{\textcircled{\scriptsize 2}}}(s-1),\qquad\lambda_1=2\ \&\ \lambda_2=1
\]
\explain{Expanding along the 3rd row: $d(s)=(s-1)\big[(s-1)(s-3)+1\big]=(s-1)(s^2-4s+4)=(s-1)(s-2)^2$; the same follows from $\det(sI-\bar A)$, since $\bar A$ is block triangular.}
\[
\underbrace{\Null(A-\lambda_1I)}_{\dim=1}+\underbrace{\Null(A-\lambda_2I)}_{\dim=1}\neq\C^3
\]
\explain{$\Rightarrow$ only two linearly indep.\ eigenvectors ($[1\ 1\ 0]\T$ and $[0\ 1\ 1]\T$) $\Rightarrow$ \key{Not diagonalizable} as the dimensions do not add up ($1+1<3$).}
\[
\underset{3\ \checkmark}{\C^3}=\underbrace{\Null\big((A-\lambda_1I)^{\text{\textcircled{\scriptsize 2}}}\big)}_{2}\oplus\underbrace{\Null(A-\lambda_2I)}_{1}
\]
% corrected: the circled exponent 2 carries a small mark (*, +) in the notes linking it to the multiplicity in d(s); dropped.
\begin{insight}
The circled exponent is the same number in both places: $\lambda_1=2$ is a double root of $d(s)$, and the right invariant subspace for it is $\Null\big((A-2I)^2\big)$ (a generalized eigenspace), not the eigenspace. Indeed
\[
(A-2I)^2=\bmat{0&0&0\\0&0&1\\0&0&1}\ \Rightarrow\ \Null\big((A-2I)^2\big)=\{x:\ x_3=0\}=\Span\{b_1^1,b_1^2\}=M_1 .
\]
By the result on p.~52 ($\Null(p(A))$ is invariant), this subspace is automatically invariant under $A$. So $A$ cannot be diagonalized, but it is block diagonal in this basis, with an upper-triangular $2\times2$ block for $\lambda_1$ (a Jordan block).
\end{insight}
\end{example}

\subsection{Distinct Eigenvalues and Diagonalization}\label{ss:diagonalization}
Let $A\in\C^{n\times n}$, $\lambda_1\neq\lambda_2$, $Ae_1=\lambda_1e_1$, $Ae_2=\lambda_2e_2$.
\begin{align*}
&\alpha_1e_1+\alpha_2e_2=0\\
\Rightarrow\ &(A-\lambda_1I)(\alpha_1e_1+\alpha_2e_2)=0\\
\Rightarrow\ &\alpha_1\underbrace{\cancel{(A-\lambda_1I)e_1}}_{0}+\alpha_2\underbrace{(A-\lambda_1I)e_2}_{\neq0\ \text{(see below)}}=0\ \Rightarrow\ \boxed{\alpha_2=0}
\end{align*}
where
\begin{align*}
(A-\lambda_1I)e_1&=Ae_1-\lambda_1e_1=\lambda_1e_1-\lambda_1e_1=0\\
(A-\lambda_1I)e_2&=Ae_2-\lambda_1e_2=\lambda_2e_2-\lambda_1e_2=(\lambda_2-\lambda_1)e_2\neq0
\end{align*}
% corrected: notes had (A - lambda_2 I)e_2 in the second term (which would be 0); applying (A - lambda_1 I) gives (A - lambda_1 I)e_2, computed below.
% corrected: notes had A in C^n
Then $\alpha_1e_1=0$ with $e_1\neq0$ gives $\alpha_1=0$ as well, so \key{$e_1,e_2$ are linearly independent}.

\begin{insight}[Eigenvectors of distinct eigenvalues are independent]
The same trick works for $k$ distinct eigenvalues: if $\sum_{i=1}^k\alpha_ie_i=0$, apply $\prod_{i\neq j}(A-\lambda_iI)$. Every term except the $j$-th is annihilated, leaving $\alpha_j\prod_{i\neq j}(\lambda_j-\lambda_i)\,e_j=0$, so $\alpha_j=0$ for every $j$. Consequently, \key{if $A\in\C^{n\times n}$ has $n$ distinct eigenvalues, it is diagonalizable}: its eigenvectors form a basis, $\C^n=\Null(A-\lambda_1I)\oplus\dots\oplus\Null(A-\lambda_nI)$, and $P^{-1}AP=\diag(\lambda_1,\dots,\lambda_n)$ with $P=[e_1\ \dots\ e_n]$. The example above shows this can fail when an eigenvalue is repeated.
\end{insight}

% ===== Page 58 =====
\begin{theorem}
\leavevmode\origpage{58}Let $A$ be an $n\times n$ matrix with $n$ distinct eigenvalues, $\lambda_1,\lambda_2,\dots,\lambda_n$ where $\lambda_i\neq\lambda_j$ when $i\neq j$. Then
\[
\begin{array}{rccccccc@{\;}l}
\C^n = & \Null(A-\lambda_1 I) & \oplus & \Null(A-\lambda_2 I) & \oplus\cdots\oplus & \Null(A-\lambda_n I) & & & \\
 & \Span\{e_1\} & & \Span\{e_2\} & & \Span\{e_n\} & & & \\[4pt]
 & 1 & + & 1 & +\cdots+ & 1 & = & n &
\end{array}
\]
\end{theorem}
Since $\Null(A-\lambda_i I)$ is invariant under $A$, \key{$A$ has a diagonal representation}:
\[
\bar A=\begin{bmatrix}\lambda_1 & & & 0\\ & \lambda_2 & & \\ & & \ddots & \\ 0 & & & \lambda_n\end{bmatrix},
\qquad \bar A = P^{-1}AP \quad\text{where } P=\big[\,e_1 \mid e_2 \mid \dots \mid e_n\,\big].
\]
% corrected: notes label the diagonal matrix "A"; it is \bar A = P^{-1}AP

\begin{proof}[Proof of the linear independence of $e_1,\ e_2,\ \dots,\ e_n$]
Suppose
\[ \alpha_1 e_1+\alpha_2 e_2+\dots+\alpha_n e_n = 0 \]
and use
\[ (A-\lambda_i I)e_j = Ae_j-\lambda_i e_j = \lambda_j e_j - \lambda_i e_j = (\lambda_j-\lambda_i)e_j . \]
Multiply with $(A-\lambda_1 I)$, then with $(A-\lambda_2 I)$, and so on:
\begin{align*}
\times(A-\lambda_1 I):&\quad 0 + (\lambda_2-\lambda_1)\alpha_2 e_2 + \alpha_3(\lambda_3-\lambda_1)e_3 + \dots + \alpha_n(\lambda_n-\lambda_1)e_n = 0\\
\times(A-\lambda_2 I):&\quad 0+0+\alpha_3(\lambda_3-\lambda_1)(\lambda_3-\lambda_2)e_3+\dots+\alpha_n(\lambda_n-\lambda_1)(\lambda_n-\lambda_2)e_n = 0\\
&\quad\ \vdots\\
&\quad \alpha_n\underbrace{(\lambda_n-\lambda_1)(\lambda_n-\lambda_2)\cdots(\lambda_n-\lambda_{n-1})}_{\neq 0}\,\underbrace{e_n}_{\neq 0} = 0\\
&\quad\Longrightarrow\quad \alpha_n = 0 \;\Longrightarrow\; \forall \alpha_i = 0\qedhere
\end{align*}
\end{proof}
\explain{The product is nonzero for distinct eigenvalues (this case), and $e_n$ is nonzero.}
\begin{insight}[Completing the argument]
After the $n-1$ multiplications only the $e_n$-term survives, so $\alpha_n=0$. The original relation then has only the terms $e_1,\dots,e_{n-1}$; repeating the procedure with the factors $(A-\lambda_1I),\dots,(A-\lambda_{n-2}I)$ gives $\alpha_{n-1}=0$, and so on; hence all $\alpha_i=0$. Thus $\{e_1,\dots,e_n\}$ is a set of $n$ linearly independent vectors in $\C^n$, i.e.\ a basis, and each $\Null(A-\lambda_iI)$ is exactly one-dimensional (it is at least $1$-dimensional, and the dimensions cannot add up to more than $n$). This proves the direct sum in the theorem.
\end{insight}

\sep
\[
P=\big[\,e_1\mid e_2\mid\dots\mid e_n\,\big], \qquad \bar A = P^{-1}AP \;\Rightarrow\; P\bar A = AP
\]
\begin{align*}
\big[\,e_1\mid e_2\mid\dots\mid e_n\,\big]
\begin{bmatrix}\lambda_1 & 0 & \cdots & 0\\ 0 & \lambda_2 & & \vdots\\ \vdots & & \ddots & \vdots\\ 0 & \cdots & \cdots & \lambda_n\end{bmatrix}
&= A\big[\,e_1\mid e_2\mid\dots\mid e_n\,\big]\\
&= \big[\,Ae_1\mid Ae_2\mid\dots\mid Ae_n\,\big]\\
&= \big[\,\lambda_1 e_1 \;\; \lambda_2 e_2 \;\; \dots\;\; \lambda_n e_n\,\big]
\end{align*}

\begin{keybox}
\noindent\hfil$\displaystyle\begin{aligned}
\therefore\ \text{Distinct eigenvalues} &\;\Rightarrow\; \text{Diagonalizable matrix}\\
&\;\Rightarrow\; \text{we can find $P$ matrix \& $\bar A$}\\[2pt]
\text{Distinct eigenvalues} &\;\not\Leftarrow\; \text{Diagonalizable}
\end{aligned}$\hfil\par\smallskip
Contrapositive ($p\Rightarrow q \;\equiv\; \bar q\Rightarrow \bar p$):\quad Non-diagonalizable $\;\Rightarrow\;$ repeated eigenvalues.
\end{keybox}
Two diagonal matrices with repeated eigenvalues:
\[
A=\bmat{1&0&0\\0&1&0\\0&0&1},\quad d(s)=(s-1)^3
\qquad\qquad
\bmat{2&0&0\\0&2&0\\0&0&1}
\]
\explain{Diagonal, but repeated eigenvalues.}
Both matrices are already diagonal (so trivially diagonalizable with $P=I$), yet their eigenvalues are repeated: this is why the converse implication fails.
